% 原书分章 PDF 第 5—7 页：第 11.2—11.4 节。
\section{因果模型与感兴趣的参数}

下面利用结构因果模型（structural causal model），定义感兴趣的因果参数。我们假设，观测纵向数据结构中的每个分量，都是一组已观测父变量和一个未测量外生误差项的函数。假定 $L_k$ 的已观测父变量为 $\bar L_{k-1}$ 和 $\bar A_k$，而 $A_k$ 的父变量为 $\bar L_k$ 和 $\bar A_{k-1}$。可以定义一个干预后分布，表示对 $A_0,\ldots,A_K$ 施加指定干预、分别将它们设为 $a_0,\ldots,a_K$ 后，数据所具有的分布。将静态干预 $\bar a$ 下真实的干预后分布记为 $P_0^{\bar a}$，并定义具有该分布的反事实随机变量 $L^{\bar a}=(L_1^{\bar a},\ldots,L_{K+1}^{\bar a})$。我们希望评价两种干预下的平均反事实结局 $E_{P_0^{\bar a}}(L_{K+1}^{\bar a})$。第一种干预分配 $A_0=1$，即接受试验疫苗；第二种干预分配 $A_0=0$，即接受对照疫苗。两者随后均分配 $A_1=\cdots=A_K=1$，即个体在整个研究期间始终处于观察之下。这些干预可能显得过于严格：只要感染发生在时点 $k$ 之前，即使个体后来失访，其在时点 $k$ 的感染状态仍是已知的。因此，对 $A_k$ 施加将其设为 $A_k=0$、而在其他情形不作干预的随机干预，可能更为适当。不过，可以证明，两种干预会得到相同的观测数据参数。所关注的两种静态干预，仅在 $A_0$ 的分配上有所不同。因此，以下简写 $a=1$ 表示分配处理且无删失的干预，$a=0$ 表示分配对照且无删失的干预。

我们希望估计疫苗组和安慰剂组中，匹配感染与不匹配感染各自的累积发生率。对 $j=1,2$ 和 $a=0,1$，定义
\[
 \psi_{0,j}^a=E_{P_0^a}(L_{K+1,j}^a).
\]
利用这些量，可以定义基因型特异的累积疫苗效力度量：
\[
 \mathrm{VE}_j=1-\frac{\psi_{0,j}^1}{\psi_{0,j}^0},\qquad j=1,2.
\]
它表示疫苗引起的第 $j$ 类感染累积发生率在乘法尺度上的降低。疫苗效力接近 1，说明疫苗高度有效；较小的正值表示疫苗具有中等效力；小于零则表示疫苗有害。我们还定义疫苗筛选效应为
\[
 \mathrm{VSE}=
 \frac{\psi_{0,1}^0/\psi_{0,1}^1}{\psi_{0,2}^0/\psi_{0,2}^1},
\]
即两种感染类型（匹配与不匹配）下，因果累积风险比（安慰剂相对于疫苗）的比值。注意，若疫苗对匹配感染的效力更高，疫苗筛选效应就大于 1；若对匹配感染的效力更低，该效应就小于 1。

为便于叙述，下面重点讨论 $\psi_{0,1}^1$ 的估计。由于感染类型和处理组的编号是任意的，同样的方法也可用于估计所关注的四个累积发生率。再将这些估计结合起来，就能估计疫苗效力和疫苗筛选效应。此外，$K$ 的选择也是任意的，因此，这些方法可用于任意一次临床访视时点的累积发生率、疫苗效力和疫苗筛选效应的逐点估计。

\section{识别}

在序贯随机化（sequential randomization）和正值性（positivity）假设下，可以利用观测数据，识别反事实变量 $L^a$ 的分布。下面的识别结论以 Bang 和 Robins（2005）的一般结果为基础。从最后一个时点 $K+1$ 开始，定义
\[
 \bar Q_{0,K+1}(\bar\ell_K)
 =\int \ell_{K+1,1}\,dQ_{0,L_{K+1}}(\bar\ell_{K+1}),
\]
其中，$Q_{0,L_{K+1}}$ 表示由 $P_0$ 所确定的、给定 $\bar A_K=\bar1_{K+1}$ 和 $\bar L_K$ 时，$L_{K+1}$ 的条件分布。给定 $\bar Q_{0,K+1}$，再定义
\[
 \bar Q_{0,K}(\bar\ell_{K-1})
 =\int \bar Q_{0,K+1}(\bar\ell_K)\,dQ_{0,L_K}(\bar\ell_K),
\]
即 $\bar Q_{0,K+1}$ 关于 $Q_{L_K,0}$ 的均值；这里，$Q_{L_K,0}$ 是由 $P_0$ 所确定的、给定 $\bar A_K=\bar1_K$ 和 $\bar L_{K-1}$ 时，$L_K$ 的条件分布。继续这一过程，对 $k=1,\ldots,K-1$，定义
\[
 \bar Q_{0,k}(\bar\ell_{k-1})
 =\int \bar Q_{0,k+1}(\bar\ell_k)\,dQ_{0,L_k}(\bar\ell_k).
\]
最后，定义
\[
 \bar Q_{0,0}=\int \bar Q_{0,1}(\ell_0)\,dQ_{0,L_0}(\ell_0).
\]
用 $\bar Q=\bar Q(P)=(\bar Q_k(P):k)$ 表示这组迭代均值，并用
$Q=Q(P)=\{\bar Q(P),Q_{L_0}(P)\}$
表示由 $P\in\mathcal M$ 所确定的迭代均值与基线协变量分布的集合。在前述因果假设下，$\bar Q_{0,0}$ 等于所关注的反事实参数。因此，在指定假设下，已经得到 $\psi_{0,1}^1=\Psi(Q_0)$，其中 $\Psi:\mathcal M\to(0,1)$ 由上述迭代均值构造定义。

\section{有效影响函数}

对 $k=1,\ldots,K+1$，定义
\[
 G_{A_k}=G_{A_k}(P)=P(A_k=1\mid\bar A_{k-1}=\bar1_k,\bar L_k),
\]
以及
\[
 H_k(G)(o)=
 \frac{I(\bar a_{k-1}=\bar1_k)}
 {\displaystyle\prod_{m=0}^{k-1}G_{A_m}(\bar\ell_k)}.
\]
在模型 $\mathcal M$ 中，$\Psi$ 在 $(Q,g)$ 处的有效影响函数（efficient influence function）为
\[
 D^*(Q,G)(o)=\sum_{k=0}^{K+1}D_k^*(Q,G)(o),
\]
其中
\[
 \begin{aligned}
 D_{K+1}^*(Q,G)(o)
 &=H_{K+1}(G)(o)\{\ell_{K+1,1}-\bar Q_{K+1}(\bar\ell_K)\},\\
 D_k^*(Q,G)(o)
 &=H_k(G)(o)\{\bar Q_{k+1}(\bar\ell_k)-\bar Q_k(\bar\ell_{k-1})\},
 \quad k=1,\ldots,K,\\
 D_0^*(Q)(o)&=\bar Q_1(\ell_0)-\Psi(\bar Q).
 \end{aligned}
\]
